\chapter{Boltzmann Transport and Berry Phase}
\label{ch:wannier-transport-berry}

Two modules on \menu{Wannier Functions} that consume a Wannier Hamiltonian
$H(\mathbf{R})$ and compute from it directly, in process --- no job is
launched and no interpreter is started.

\begin{calnote}
Both are native Calango implementations. Wannier90, postw90 and BoltzWann are
references for the formulas, conventions and feature scope; they are not
dependencies. No external binary is invoked, nothing is linked against them,
and none of their auxiliary output files is required. Reading an
\file{\_hr.dat} you already have is done by Calango's own parser --- the file
is a plain text table of hopping integrals.
\end{calnote}

Both take their input from the same row, which offers three sources. The
normal one is \emph{From a completed run}: a Wannier Functions run finished in
this session, whose \file{wannier\_hr.dat}, cell, centres and spreads are read
straight from its directory, so nothing outside Calango is involved. The others
are opening an \file{\_hr.dat} you already have, and loading one of two
built-in demonstration models (a simple-cubic metal, and the two-band
Qi--Wu--Zhang Chern insulator). Those are the models the unit tests use, so
what the panel reports can be checked against the test output directly.

A run started before Calango wrote $H(\mathbf{R})$ says so when picked, and
says to re-run the Wannierization on the same baseline --- distinguished from a
run that recorded the file but no longer has it, which needs a different fix.

\section{Boltzmann Transport}

\menu{Wannier Functions}\menusep\menu{Boltzmann Transport\ldots}

Electronic and thermoelectric transport in the constant relaxation-time
approximation. Everything derives from the transport distribution function
\[
  \Sigma_{\alpha\beta}(\varepsilon) = \frac{1}{V N_k}\sum_{n,\mathbf{k}}
    \tau\, v_{n\alpha} v_{n\beta}\,
    \delta(\varepsilon - \varepsilon_{n\mathbf{k}}),
\]
whose Onsager moments
$L^{(m)} = \int \! d\varepsilon\, \Sigma\, (\varepsilon-\mu)^m
(-\partial f/\partial\varepsilon)$ give the conductivity
$\sigma = e^2 L^{(0)}$, the Seebeck tensor
$S = (eT)^{-1}(L^{(0)})^{-1}L^{(1)}$, and the electronic thermal
conductivity
\[
  \kappa_e = \frac{1}{e^2 T}\left[L^{(2)}
    - L^{(1)}\bigl(L^{(0)}\bigr)^{-1}L^{(1)}\right].
\]
That is the \emph{zero-current} thermal conductivity, which is the one $zT$
wants. The subtracted term is the Peltier heat carried by the compensating
current; dropping it inflates $\kappa_e$ and deflates $zT$. The panel also
reports the power factor $S^2\sigma$, the figure of merit
$zT = S^2\sigma T/(\kappa_e + \kappa_L)$, and the carrier concentration.

\subsection{Band velocities}

Velocities come from $\partial H/\partial \mathbf{k}$ in the Wannier gauge,
taken analytically. This is not an optimisation: at a band crossing the
eigenvalue branches swap, so a finite difference of \emph{sorted} eigenvalues
reports a velocity that jumps discontinuously and can have the wrong sign
entirely. The matrix gradient has no such problem, because the degeneracy
lives in the eigenvectors.

\subsection{The relaxation time}

$\tau$ enters $\sigma$ and $\kappa_e$ linearly and cancels \emph{exactly} in
$S$, which is a ratio of two moments that each carry one factor of it. A
constant-$\tau$ Seebeck coefficient is therefore a genuine prediction, while
$\sigma$ and $\kappa_e$ are only as good as the $\tau$ supplied. The lattice
thermal conductivity $\kappa_L$ is a phonon quantity that nothing in an
electronic structure determines, so it is a user input and is needed only for
$zT$.

\begin{calwarn}
The smearing that resolves the delta function has to sit between two bounds:
\emph{above} the mean level spacing of the k-mesh, or $\Sigma(\varepsilon)$
breaks into isolated spikes, and \emph{below} $k_BT$, or the smearing acts as
extra temperature and the Wiedemann--Franz ratio comes out low. At 300\,K,
$k_BT$ is 26\,meV. The panel reports the measured level spacing against both
after every run.
\end{calwarn}

\section{Berry Phase}

\menu{Wannier Functions}\menusep\menu{Berry Phase\ldots}

\begin{center}
\small
\begin{tabular}{@{}lp{7.6cm}@{}}
\toprule
Quantity & How it is obtained \\
\midrule
Berry phase on a closed path & Wilson loop: product of overlap matrices \\
Berry curvature $\Omega(\mathbf{k})$ & Kubo sum over band pairs of
  $\partial H/\partial k$ elements \\
Curvature map & the same over a k-plane, drawn with the app's colormaps \\
Anomalous Hall conductivity & BZ integral of $\Omega$, in SI and in $e^2/h$ \\
Electric polarization & Berry phase (modern theory), with its quantum \\
Hybrid Wannier centres & Wilson-loop eigenphases vs transverse $k$ \\
\bottomrule
\end{tabular}
\end{center}

\subsection{Gauge}

This is the whole difficulty of the subject, so the conventions are fixed
explicitly. Berry phases are Wilson loops: an arbitrary phase on an
eigenvector enters once as a bra and once as a ket and cancels around the
loop, so the result does not depend on what the eigensolver returned. Berry
curvature uses the Kubo form, in which every factor is a matrix element
between eigenstates, so the phases cancel pairwise and no eigenvector is ever
differentiated. $H(\mathbf{k})$ is built in convention~I --- orbital positions
\emph{not} in the phase --- so $H(\mathbf{k}+\mathbf{G}) = H(\mathbf{k})$
exactly and a loop closes with no correction factor.

\subsection{Relation to Topological Invariants}

The hybrid Wannier centre flow is the same object
\menu{Wannier Functions}\menusep\menu{Topological Invariants\ldots} plots.
That feature generates a Python script and runs it against a completed
Wannierization; this one evaluates the same quantity in process from
$H(\mathbf{R})$. The two now overlap more than they used to: on a
\textbf{GPAW}-sourced run the generated script reopens the \file{.gpw} and
transports the Bloch states with \texttt{gpaw.berryphase}, but on a
\textbf{VASP}-sourced one it builds a Wilson loop from the same
\file{wannier\_hr.dat} this panel reads --- so on that route they are the same
algorithm reached two ways, and on the GPAW route they remain genuinely
independent. Either way this panel reports the Chern number twice, once from
the curvature integral and once from the centre winding, so that a
disagreement flags under-sampled $k$.

One thing the $H(\mathbf{R})$ route cannot do, on either side: a wannier90
\file{\_hr.dat} is an $N_w \times N_w$ matrix with \emph{no spin labelling}.
For a spinor (\texttt{LSORBIT}) run its Wannier functions are spinors and
$N_w$ counts them, but the file records neither the up/down decomposition nor
the convention used to order it. Neither invariant reads a spin expectation
--- the Chern winding and the Soluyanov--Vanderbilt largest-gap $Z_2$
are computed from the centres alone --- so this costs a plot series, not a
result, and \file{topology.json} carries \texttt{spin: null} rather than an
invented one.

\section{Validation}

Both modules are covered by unit tests against results known independently of
the code: the Wiedemann--Franz limit
$\kappa_e/\sigma T \to \pi^2 k_B^2/3e^2 = 2.44\times10^{-8}$\,W$\Omega$/K$^2$,
the Seebeck sign for electron and hole doping, the exact $\tau$-scaling
identities, the Qi--Wu--Zhang Chern number ($\pm1$ inside $|m|<2$, zero
outside), the Hall conductance quantised at $e^2/h$, and the SSH Zak phase,
which comes out exactly $0$ and $\pi$ in the two phases.
